\section{Introduction}

A join between a large fact table and a smaller dimension table is the most
common heavy operation in an analytics pipeline, and skew in the join key is
the most common reason such a join is slow. When one key carries a large share
of the rows, hash partitioning sends them all to a single reducer, one task
runs far longer than its peers, and the stage finishes when that task finishes.
Adding compute does not help: the straggler is serial. The effect was
characterised for parallel joins three decades ago~\cite{walton1991taxonomy,
dewitt1992skew} and re-confirmed at production scale in every generation of
big-data system since~\cite{chaiken2008scope, ananthanarayanan2010mantri,
kwon2012skewtune}.

Apache Spark~\cite{zaharia2012rdd, armbrust2015sparksql} addresses skew
automatically. Since version 3.0, the \emph{Adaptive Query Execution} (AQE)
framework re-optimises a plan between stages using statistics observed at
runtime, and its \texttt{OptimizeSkewedJoin} rule detects a shuffle partition
that is disproportionately large, splits it, and replicates the matching
partition on the other side so that every split finds its
partners~\cite{databricks_aqe, spark_aqe_docs, xue2024adaptive}.

And yet practitioners keep salting by hand. \emph{Salting} appends a random
suffix drawn from $\{0,\dots,k-1\}$ to the join key on the skewed side and
replicates the matching rows $k$ times on the other side, so that one heavy key
becomes $k$ lighter ones. The technique is ubiquitous in industry and, as we
establish in Section~\ref{sec:background}, entirely absent from the
peer-reviewed literature under that name. It is transmitted through vendor
documentation, conference talks and blog posts, none of which prices what it
costs.

That silence matters, because salting is not free. Every additional salt value
multiplies the rows that must be replicated and shuffled on the opposite side.
The straggler term falls as $1/k$; the replication term rises linearly in $k$.
Total time is therefore \emph{convex} in $k$ with an interior minimum, and the
standard advice --- ``pick a reasonable $k$, say 16 or 32'' --- is not advice at
all. It names a point on a curve whose shape nobody has published.

\subsection{The questions this paper answers}

\begin{enumerate}
\item Given a join's key distribution and Spark's configuration, will AQE's
skew-join rule suffice on its own, or is explicit salting still required?
\item If salting is required, what value of $k$ minimises total cost, and can
it be derived from statistics available before the query runs rather than found
by trial and error?
\item Does combining AQE with salting help, hurt, or do nothing?
\end{enumerate}

\subsection{Contributions}

\begin{itemize}
\item \textbf{A cost model for salt cardinality.} We formulate the salting
trade-off as $T(k) = aH/k + bP(k-1) + c$, where $H$ and $P$ are the hot-key row
counts on the skewed and probe sides, and derive the closed-form optimum
$k^{*}=\sqrt{aH/(bP)}$. The single free ratio $\gamma = a/b$ is a property of
the engine and hardware, calibrated once and reused
(Section~\ref{sec:costmodel}).

\item \textbf{Two ceilings that bound the useful range.} Beyond
$k=\rho$, where $\rho$ is the hot key's rows divided by the mean rows per
partition, the hot task is no longer the straggler and further splitting
relieves nothing; beyond $k=C$, the number of concurrently schedulable task
slots, the sub-partitions queue rather than overlap. \emph{Salting beyond
$\min(\rho, C)$ is always waste} --- a bound that follows directly from the
model and that we could not find stated anywhere in the practitioner
literature.

\item \textbf{A decision procedure.} A rule that predicts, from data statistics
and Spark configuration alone, whether AQE will trigger, whether it will
suffice, and what $k$ to use otherwise. Its hinge is the distinction between
one-sided skew, which AQE is designed for, and two-sided skew, in which
split-and-replicate degenerates because each split inherits a large partner
(Section~\ref{sec:decision}).

\item \textbf{\textsc{skewbench}, an open harness.} A reproducible benchmark
that generates Zipfian-skewed workloads with an aerospace engine-telemetry
schema, executes six join strategies, and recovers hardware-independent metrics
--- shuffle bytes, spill bytes, and the within-stage ratio of maximum to median
task duration --- by parsing the Spark event log. It runs on a commodity laptop
(Section~\ref{sec:methodology}).

\item \textbf{An empirical evaluation} over 48 configurations, which confirms
the decay of AQE's benefit as the probe side thickens, identifies the
precondition that actually blocked Spark's skew-join rule from ever firing, and
quantifies the replication-cost gap between the two salting variants. It also
reports, as a negative result, that the convexity of $T(k)$ does not resolve at
single-node scale for selective salting (Section~\ref{sec:results}).
\end{itemize}

\subsection{Scope, stated plainly}

This is an evaluation and decision-procedure paper, not an algorithm paper. We
propose no new partitioner; the space of skew-resilient partitioning algorithms
is already well populated~\cite{irandoost2019skewslr, gufler2011skew,
myung2016handling, shen2020srspark, tang2021intermediate}. Our claim is that
the prior question --- \emph{which of the mechanisms that already ship should a
practitioner use, and how should it be parameterised} --- has gone unanswered,
and that it is answerable.

All measurements were taken on a single node. We are explicit about what that
does and does not license. Absolute wall-clock times do not transfer to a
cluster and we make no claim that they do; shuffle volume, spill volume and
within-stage task-time ratios are properties of the physical plan and the data
partitioning rather than of the machine, and it is on those that the
quantitative claims rest. Section~\ref{sec:threats} treats this at length.
